\documentclass[11pt]{article}
\usepackage[a4paper,margin=25mm]{geometry}
\usepackage{amsmath,amssymb,amsthm,hyperref}
\newtheorem{theorem}{Theorem}\newtheorem{lemma}{Lemma}
\newcommand{\PP}{\mathcal P}
\title{Two fair pure densities with polynomial-denominator cell means}
\author{Complete proof; three independent reviews passed}\date{30 September 2026}
\begin{document}\maketitle
Here and below ``polynomial'' means bounded by one fixed power of $m+1$.
Constants and exponents are absolute; their optimization is not attempted.
This note strengthens the rational two-density completion by controlling
one common denominator for all cell masses and first moments.

\begin{theorem}\label{main}
For every $m\ge2$ there are rational piecewise polynomial probability
densities $f_A,f_B$ on $[0,1]$ such that $A,B$ together with $m$
independent uniforms are fully permutation-fair. There is a partition
into polynomially many rational cells with these additional properties:
\begin{enumerate}
\item On each open cell exactly one of $f_A,f_B$ is positive. Its
conditional density, affinely rescaled to $[0,1]$, lies between two
absolute positive constants and has polynomial degree.
\item Every cell length and its positive mass are at least an inverse
polynomial.
\item All endpoints, all cell masses, and all unnormalized first cell
moments have one common polynomial rational denominator.
\end{enumerate}
No denominator bound is asserted for the density coefficients or higher
cell moments.
\end{theorem}

\begin{lemma}[Conditioning cumulative mass and first moment]\label{conditioning}
Let $J\ge2$, $d>2(J+4)+1$ with $d=O(J)$, and put
\[
 H(t)=\int_0^tP_d(2u-1)du,\quad
 R_0(t)=t-\eta H(t),\quad
 R_1(t)=\int_0^t u[1-\eta P_d(2u-1)]du,
\]
where $0<\eta\le1/2$. The family
$\{\phi_j, R_0\phi_j,R_1\phi_j:0\le j<J\}$, where $\phi_j$
are orthonormal shifted Legendre polynomials, has Gram matrix bounded
below by $\eta^2(J+1)^{-C}$ for an absolute $C$.
The same holds on separated midpoint-like exact quadrature nodes of
sufficient degree, and under sufficiently small polynomial perturbations.
\end{lemma}
\begin{proof}
Write $w=t(1-t)$ and $\lambda=d(d+1)$. The primitive identity
\[
 \int_0^tH(u)du=\frac{w'H-wH'}{\lambda-2}
\]
gives
\[
 R_1=\frac{t^2}2-\eta\left(t-\frac{w'}{\lambda-2}\right)H
                    -\frac{\eta w}{\lambda-2}H'.
\]
Thus $p+qR_0+vR_1$ is $\widetilde p+H\widetilde q+H'\widetilde v$,
with $\widetilde v=-\eta wv/(\lambda-2)$, and all coefficient degrees
at most $J+2$. The three-band lemma in
\path{../rational_designs/asymptotic/ordered_prefix_legendre_tools.tex}
recovers these coefficients with a polynomial norm loss. Multiplication
by $w$ on $\PP_{J-1}$ has inverse norm at most $C(J+1)^2$: delete
endpoint intervals of length $c(J+1)^{-2}$ and apply Nikolskii to
retain half the squared norm. Recover $v$, then $q$, then $p$.
The total inverse norm is at most $\eta^{-1}(J+1)^C$.
An exact quadrature for the squares reproduces this Gram matrix.
Polynomial evaluation and derivative bounds preserve it under sufficiently
small perturbations.
\end{proof}

\begin{proof}[Proof of Theorem~\ref{main}]
Set $J=m+3$, choose $d$ as in Lemma~\ref{conditioning}, and choose
$\eta=(m+1)^{-A}$ with $A$ a sufficiently large fixed integer to be
specified by the positivity estimates below. Take a separated real
equal-weight rule $b_q^0$, $1\le q\le K$, of degree at least
$2(d+J+3)$, with $K=C(m+1)^4$ rounded to an integer and midpoint
discrepancy $O(1/K)$. Such a rule, with endpoint and consecutive gaps
at least $c/K$, is provided by the midpoint-grid Newton lemma in
\path{../rational_designs/asymptotic/two_discrete_step_compatibility.tex}.
The exact degree can be increased by an absolute factor without changing
this bound on $K$.

Legendre orthogonality and Fubini give, for $0\le j\le J$,
\begin{equation}\label{baseline}
 \frac1K\sum_q(b_q^0)^j=\frac1{j+1},\quad
 \frac1K\sum_qR_0(b_q^0)(b_q^0)^j=\frac1{j+2},\quad
 \frac1K\sum_qR_1(b_q^0)(b_q^0)^j=\frac1{2(j+3)}.
\end{equation}
For the third identity the correction integral is
$-(\eta/(j+1))\int_0^1uP_d(2u-1)(1-u^{j+1})du=0$.
In particular the two feature means are $1/2$ and $1/6$.

Choose one polynomial integer $L$, divisible by $6$, with sufficiently
large fixed exponent. Independently round the arrays $R_0(b_q^0)$ and
$R_1(b_q^0)$ to $r_{0q},r_{1q}\in L^{-1}\mathbb Z$, preserving their
sums $K/2$ and $K/6$ exactly. Balanced rounding gives error at most $1/L$.
For the remaining equations in~\eqref{baseline}, keep $r_0,r_1$ fixed
and vary the $b_q$. Use antiderivatives of orthonormal polynomials as
the nonconstant moment basis. The derivative matrix has rows given by
$K^{-1}\phi_j(b_q)$, $K^{-1}r_{0q}\phi_j(b_q)$, and
$K^{-1}r_{1q}\phi_j(b_q)$, $j<J$.
Lemma~\ref{conditioning} provides a right inverse of polynomial norm
because $\eta$ is an inverse polynomial. The rounding residual and
Jacobian perturbation are at most a polynomial times $1/L$.
Polynomial derivative bounds give a fixed-inverse contraction in a
ball of radius $\operatorname{poly}(m)/L$. Taking the exponent in $L$
sufficiently large gives real separated nodes $b_q$ satisfying
\begin{equation}\label{ranks}
 K^{-1}\sum_q b_q^j=\frac1{j+1},\quad
 K^{-1}\sum_q r_{0q}b_q^j=\frac1{j+2},\quad
 K^{-1}\sum_q r_{1q}b_q^j=\frac1{2(j+3)}
 \qquad(0\le j\le J),
\end{equation}
with $\max_q|b_q-b_q^0|\le\operatorname{poly}(m)/L$.
The empirical feature Gram bound survives. In particular
\[
 |r_{0q}-b_q|+|r_{1q}-b_q^2/2|\le C\eta
\]
after increasing $L$. The array $r_0$ is strictly increasing and
interior, since its baseline has derivative $1-\eta P_d\ge1/2$.

Apply the mean-preserving rational box theorem in
\path{../rational_designs/asymptotic/polynomial_mean_box_smoothing.tex}
with the three fixed features $(1,r_0,r_1)$ and degree $J$.
Its low-degree arithmetic targets are the fixed rational numbers
$1/2,1/3,1/8$, so its hypotheses hold. Obtain rational boxes
$J_q=[c_q-\epsilon,c_q+\epsilon]$, each of mass $1/K$, with positive
quadratic conditional densities. Their endpoints and means $c_q$ lie
on one polynomial rational grid. All three families in~\eqref{ranks}
remain exact with box integrals in place of node evaluations.
The exponent in the rational halfwidth $\epsilon$ can be chosen
arbitrarily large in advance, and the center displacement can be made
smaller than any prescribed inverse polynomial. The complementary gaps
still have lengths at least $\Delta=c/K$.
Let this smoothed density be $f_B=\beta$.

On the complementary gaps prescribe cumulative mass $r_{0q}$ and
cumulative first moment $r_{1q}$ before box $q$. At the two outer
endpoints prescribe $(0,0)$ and $(1,1/2)$. Let $g$ be the unique affine
density on each gap with the resulting mass and first moment, and zero
on the boxes. All its coefficients are rational. The elementary
rescaled two-moment system gives
\[
 \|g-1\|_{L^\infty(G)}
 \le C\Delta^{-2}(\eta+\epsilon+\max_q|c_q-b_q|),
\]
so it lies between $1/2$ and $3/2$ when the exponents are large enough.
The weighted identities~\eqref{ranks} give precisely compatibility on
$\mathcal S_1$ in
\path{quantitative_smoothed_quotient.tex}. That note gives
\[
 \|1-g\|_2\le
 C\Delta^{-2}(\eta+\epsilon+\max|c-b|)+\sqrt{2K\epsilon},
\]
and its completion applies if this is at most
$c(m+1)^{-7}\Delta^3$ and $\epsilon\le c\Delta/(m+1)^3$.
Choose the fixed exponent $A$ first sufficiently large, then the box
width exponent, then the rounding precisions sufficiently large.
All choices remain polynomial, and these inequalities follow.
We obtain a rational piecewise polynomial density $f_A$, between
$1/4$ and $7/4$ on the gaps, with exactly the prescribed mass and first
moment on every gap, and
\[
 \int f_A(a)F(a)da=\int_0^1F(a)da
 \quad\left(F\in\PP_m+
       \operatorname{span}\{a^rS_s(a):r+s\le m\}\right).
\]
Consequently all ordered polynomial moments of total degree at most $m$
of $(A,B)$ agree with those of two independent uniforms. For example
\[
 \mathbb E[A^rB^s\mathbf1_{A<B}]
 =\int f_A(a)a^rS_s(a)da
 =\frac1{r+1}\int b^{r+s+1}\beta(b)db
 =\frac1{(r+1)(r+s+2)}.
\]
The other chamber follows from the marginal moments and independence.
Integrating the $m$ uniform variables in a fixed full order leaves a
polynomial of total degree $m$ on either chamber, proving full fairness.

Finally, every $A$-gap mass and first moment is a difference of entries
of $r_0,r_1$, hence has denominator dividing $2L$. Every $B$-box mass
is $1/K$ and its first moment is $c_q/K$. The endpoints and means
$c_q$ have one polynomial common denominator by the smoothing theorem.
Multiplying these finitely many \emph{common} denominators, rather than
individual cell denominators, yields the asserted single polynomial
grid. The gap and box lengths are inverse polynomials. The density
bounds imply the same for every positive cell mass and give the claimed
conditional density bounds after rescaling. The piece degrees are at
most $m$ on the gaps and two on the boxes.
\end{proof}
\end{document}
